%%
%% absfigure.tex
%% Documentation for the absfigure package
%%
%% Copyright (c) 2026 Christian Holz
%% Released under the MIT License. See LICENSE for details.
%%

\documentclass[a4paper]{article}

\usepackage{xcolor}
\usepackage{hyperref}

\definecolor{notes}{rgb}{.75, .3, .3}%
\definecolor{keycolor}{rgb}{.0, .3, .6}%
\definecolor{examplebg}{gray}{.93}%

\hypersetup{%
  colorlinks=true,
  linkcolor=notes,
  urlcolor=keycolor
}

\newcommand\key[1]{\textcolor{keycolor}{\texttt{#1}}}
\newcommand\pkg[1]{\textsf{#1}}
\newcommand\env[1]{\texttt{#1}}
\newcommand\cmd[1]{\texttt{\string#1}}
\newcommand\meta[1]{\ensuremath{\langle}\textit{#1}\ensuremath{\rangle}}

\makeatletter
\newenvironment{warningblock}
  {\list{}{}%
   \item[{\sffamily\bfseries\color{notes}WARNING}]\ignorespaces}
  {\endlist}
\newenvironment{noteblock}
  {\list{}{}%
   \item[{\sffamily\bfseries\color{notes}NOTE}]\ignorespaces}
  {\endlist}
\newsavebox{\absf@exbox}
\newenvironment{exampleblock}
  {\par\medskip
   \noindent{\sffamily\bfseries\color{keycolor}EXAMPLE}\par\nobreak\smallskip
   \begin{lrbox}{\absf@exbox}%
   \begin{minipage}{\dimexpr\linewidth-16pt\relax}}
  {\end{minipage}%
   \end{lrbox}%
   \begingroup\setlength{\fboxsep}{8pt}%
   \noindent\colorbox{examplebg}{\usebox{\absf@exbox}}%
   \endgroup\par\medskip}
\makeatother

\addtolength{\topmargin}{-2pc}
\addtolength{\textwidth}{4pc}
\addtolength{\oddsidemargin}{-2pc}
\addtolength{\textheight}{5pc}

\title{Absolute figure and table placement\\with the\\\pkg{absfigure} package}

\author{Christian Holz}

\date{Version 0.1\\2026/02/25}

\begin{document}

\maketitle
\tableofcontents

\bigskip

\begin{noteblock}
  This is version~0.1 of \pkg{absfigure}. The package is functional
  but should be considered experimental. Please report any issues.
\end{noteblock}

%% ============================================================
\section{Introduction}
%% ============================================================

\LaTeX's float mechanism (\verb|\begin{figure}[t/b/h/p]|) provides
only \emph{relative} placement hints. There is no built-in way to say
``put this figure at the bottom of page~6, column~2.'' The
\pkg{absfigure} package provides exactly that capability for both
figures and tables through a simple key-value interface.

Floats placed with \pkg{absfigure} integrate with \LaTeX's page
builder: text correctly reflows around them, \cmd{\caption},
\cmd{\label}, and \cmd{\ref} all work normally, and ordinary floats
continue to function alongside absolute ones. Both
\cmd{\listoffigures} and \cmd{\listoftables} pick up the placed
floats automatically.

\subsection{Requirements}

\begin{itemize}
\item \LaTeXe{} (1995/12/01 or later)
\item The \pkg{keyval} package (part of the \pkg{graphics} bundle)
\item The \pkg{environ} package (available on CTAN)
\end{itemize}

\subsection{Loading the package}

\begin{verbatim}
\usepackage{absfigure}
\end{verbatim}

No package options are currently defined.

%% ============================================================
\section{User interface}
%% ============================================================

\subsection{The \env{absfigure} environment}

\begin{quote}
\cmd{\begin}\verb|{absfigure}{|\meta{keys}\verb|}|\\
\quad\meta{figure body}\\
\cmd{\end}\verb|{absfigure}|
\end{quote}

Places a single-column figure on a specific page and column. The
\meta{figure body} may contain \cmd{\includegraphics},
\cmd{\caption}, \cmd{\label}, and any other content you would
normally put inside a \LaTeX{} \env{figure} environment.

\begin{exampleblock}
\begin{verbatim}
\begin{absfigure}{page=6, pos=b, col=2}
  \includegraphics[width=\columnwidth]{fig.pdf}
  \caption{My figure}
  \label{fig:myfig}
\end{absfigure}
\end{verbatim}
\end{exampleblock}

\subsection{The \env{absfigure*} environment}

\begin{quote}
\cmd{\begin}\verb|{absfigure*}{|\meta{keys}\verb|}|\\
\quad\meta{figure body}\\
\cmd{\end}\verb|{absfigure*}|
\end{quote}

Places a spanning (double-column) figure, analogous to \LaTeX's
\env{figure*}. The \key{col} key is ignored for this environment
since the figure spans both columns.

\begin{exampleblock}
\begin{verbatim}
\begin{absfigure*}{page=3, pos=t}
  \includegraphics[width=\textwidth]{wide.pdf}
  \caption{A wide figure}
\end{absfigure*}
\end{verbatim}
\end{exampleblock}

\subsection{The \env{abstable} environment}

\begin{quote}
\cmd{\begin}\verb|{abstable}{|\meta{keys}\verb|}|\\
\quad\meta{table body}\\
\cmd{\end}\verb|{abstable}|
\end{quote}

Places a single-column table on a specific page and column. The
\meta{table body} may contain \env{tabular}, \cmd{\caption},
\cmd{\label}, and any other content you would normally put inside a
\LaTeX{} \env{table} environment. The same keys (\key{page},
\key{pos}, \key{col}) are accepted.

\begin{exampleblock}
\begin{verbatim}
\begin{abstable}{page=4, pos=t, col=1}
  \centering
  \caption{My table}
  \label{tab:mytab}
  \begin{tabular}{lcc}
    \hline
    Item  & Value & Unit \\
    \hline
    Alpha & 1.0   & m/s  \\
    \hline
  \end{tabular}
\end{abstable}
\end{verbatim}
\end{exampleblock}

\subsection{The \env{abstable*} environment}

\begin{quote}
\cmd{\begin}\verb|{abstable*}{|\meta{keys}\verb|}|\\
\quad\meta{table body}\\
\cmd{\end}\verb|{abstable*}|
\end{quote}

Places a spanning (double-column) table, analogous to \LaTeX's
\env{table*}. The \key{col} key is ignored for this environment
since the table spans both columns.

\begin{exampleblock}
\begin{verbatim}
\begin{abstable*}{page=5, pos=b}
  \centering
  \caption{A wide table}
  \begin{tabular}{lccccc}
    \hline
    Exp & T1 & T2 & T3 & Mean & Std \\
    \hline
    A   & 4.2 & 4.5 & 4.1 & 4.27 & 0.21 \\
    \hline
  \end{tabular}
\end{abstable*}
\end{verbatim}
\end{exampleblock}

%% ============================================================
\section{Keys}
%% ============================================================

\subsection{\key{page}}

\begin{quote}
\key{page=}\meta{integer} --- Target page number (positive integer).\\
\key{page=current} --- Place on the page currently being built (default).\\
\key{page=end} --- Place on an extra page appended after the document.
\end{quote}

If \key{page} is omitted, it defaults to \key{current} --- the float
is placed on the page where the environment appears in the source.

An error is raised if \key{page} is zero, negative, or non-numeric
(other than the special values \key{current} and \key{end}).

If the target page number exceeds the number of pages in the document,
the float is never placed and a hard error is raised at the end of the
\LaTeX{} run.

\begin{warningblock}
  When using an explicit page number, the float must be defined in the
  source \emph{before} its target page is shipped out. In practice,
  group such floats right after \verb|\begin{document}|. With
  \key{page=current} (or the default), the float is placed on the
  page being built where the environment appears, so ordering is not
  a concern.
\end{warningblock}

\subsection{\key{pos}}

\begin{quote}
\key{pos=t} --- Place at the top of the page (default).\\
\key{pos=b} --- Place at the bottom of the page.
\end{quote}

The aliases \key{pos=top} and \key{pos=bottom} are also accepted.

An error is raised if \key{pos} is set to anything other than
\key{t}/\key{top} or~\key{b}/\key{bottom}.

\begin{noteblock}
  The \LaTeX{} kernel does not natively support bottom-placed spanning
  floats (there is no \verb|\@dblbotlist|). The \pkg{absfigure} package
  implements its own spanning-bottom mechanism, so \key{pos=b} works
  with all four environments (\env{absfigure}, \env{absfigure*},
  \env{abstable}, \env{abstable*}).
\end{noteblock}

\subsection{\key{col}}

\begin{quote}
\key{col=1} --- Place in the first (left) column (default).\\
\key{col=2} --- Place in the second (right) column.
\end{quote}

The aliases \key{col=l}/\key{col=left} and \key{col=r}/\key{col=right}
are also accepted.

This key is only meaningful in \verb|\twocolumn| mode. In single-column
documents, \key{col=2} (or \key{col=right}) raises an error.

For starred environments (\env{absfigure*}, \env{abstable*}), this key
is silently ignored (the float always spans both columns).

An error is raised if \key{col} is set to anything other than
\key{1}/\key{left} or~\key{2}/\key{right}.

%% ============================================================
\section{Captions, labels, and cross-references}
%% ============================================================

\cmd{\caption} and \cmd{\label} work normally inside all four
environments. They are executed at \emph{placement} time (during the
output routine), so:

\begin{itemize}
\item The figure/table counter is incremented in the order floats are
  actually placed, not the order they appear in the source.
\item \cmd{\ref} and \cmd{\pageref} resolve correctly after two
  compilations (standard \LaTeX{} behaviour for cross-references).
\item Figures appear in \cmd{\listoffigures} and tables in
  \cmd{\listoftables} automatically.
\end{itemize}

%% ============================================================
\section{Error handling}
%% ============================================================

The package performs validation on all keys and issues errors or
warnings as appropriate:

\begin{description}
\item[\key{page} omitted] %
  Defaults to \key{current} (the page being built). No error.

\item[\key{page} is zero, negative, or non-numeric] %
  \textbf{Error:} \texttt{Invalid page value `\meta{value}'}.
  The \key{page} key requires a positive integer, \key{current},
  or \key{end}.

\item[\key{page} exceeds total pages] %
  \textbf{Error} at end of run:
  \texttt{Float N (figure/table, target page=X) was never placed}.
  Compilation is interrupted; no float is silently dropped.

\item[\key{page=current}] %
  Valid. Resolves to the page number being built at definition time.

\item[\key{page=end}] %
  Valid. An extra page is appended after the document's last page,
  and the float is placed there.

\item[\key{pos} is not \key{t}/\key{top} or \key{b}/\key{bottom}] %
  \textbf{Error:} \texttt{Invalid pos value `\meta{value}'}.

\item[\key{col} is not \key{1}/\key{l}/\key{left} or \key{2}/\key{r}/\key{right}] %
  \textbf{Error:} \texttt{Invalid col value `\meta{value}'}.

\item[\key{col=2} or \key{col=right} in single-column mode] %
  \textbf{Error:} \texttt{col=2 (right) used in single-column mode}.
\end{description}

%% ============================================================
\section{How it works}
%% ============================================================

Understanding the internal mechanism may be helpful for advanced users
and for debugging.

\subsection{Core mechanism}

\LaTeX{} calls \verb|\@startcolumn| at the start of every column (and
page). The package patches this command to inject absolute figures
\emph{before} deferred floats are processed:

\begin{verbatim}
\@startcolumn (patched):
  1. \global\@colroom\@colht      % reset height
  2. detect new column vs re-entry
  3. \absf@inject@figures          % OUR HOOK
  4. \@colroom -= \absf@colreduce  % apply reduction
  5. \@tryfcolumn / deferred       % original logic
\end{verbatim}

\subsection{Re-entry handling}

The kernel may call \verb|\@startcolumn| multiple times for the same
column (via \verb|\@whilesw\if@fcolmade|). Each re-entry resets
\verb|\@colroom|, undoing any reductions. The package tracks the
current column identity (page number + column flag) and stores the
cumulative space reduction in a dimension register
(\verb|\absf@colreduce|). On re-entry, injection is skipped but the
stored reduction is re-applied.

\subsection{Injection logic}

For each pending float, the package checks whether the current
page counter (\verb|\c@page|) matches the target page and whether
\verb|\if@firstcolumn| matches the target column. On a match:

\begin{enumerate}
\item A float box register is allocated from \verb|\@freelist|.
\item The captured body tokens are typeset inside it (with
  \verb|\@captype| set to \texttt{figure} or \texttt{table}
  so \cmd{\caption}/\cmd{\label} work and counters increment
  correctly).
\item \verb|\count\@currbox| is set to encode the float type and
  placement using the kernel's \verb|\ftype@|\meta{type} multiplier.
\item The box is added to \verb|\@toplist| or \verb|\@botlist| and
  the space reduction is accumulated in \verb|\absf@colreduce|.
\end{enumerate}

For spanning floats (\env{absfigure*}, \env{abstable*}), the box is
added to \verb|\@dbltoplist| and \verb|\@colht| is reduced instead
(this reduction naturally survives re-entries since
\verb|\@startcolumn| resets \verb|\@colroom| from \verb|\@colht|).

\subsection{Why text reflows correctly}

Reducing \verb|\@colroom| inside \verb|\@startcolumn| propagates to
\verb|\vsize|, which the \TeX{} page builder uses to determine where
to break the column. Less available height means less text on the
page, leaving room for the figure.

%% ============================================================
\section{Limitations and known issues}
%% ============================================================

\begin{itemize}
\item Floats must be defined in the source \emph{before} their
  target page. The recommended practice is to group all
  \env{absfigure}/\env{abstable} environments right after
  \verb|\begin{document}|.

\item If many absolute floats target the same page, the available
  text height may become negative and \TeX{} will produce an
  overfull vbox. Plan float sizes accordingly.

\item Float numbering follows placement order, not source order.
  If you define figure~A (targeting page~5) before figure~B
  (targeting page~2), figure~B will get a lower number. The same
  applies to tables.

\item Cross-references (\cmd{\ref}, \cmd{\pageref}) require two
  compilations to resolve, as is standard for \LaTeX.

\item The package has not been tested with the \pkg{multicol}
  package. It is designed for \LaTeX's built-in \verb|\twocolumn|
  mechanism.
\end{itemize}

%% ============================================================
\section{Complete example}
%% ============================================================

\begin{verbatim}
\documentclass[twocolumn]{article}
\usepackage{graphicx}
\usepackage{lipsum}
\usepackage{absfigure}

\begin{document}

% Define figures up front
\begin{absfigure}{page=2, pos=t, col=1}
  \centering
  \includegraphics[width=\columnwidth]{fig1.pdf}
  \caption{Top of page 2, column 1}
  \label{fig:one}
\end{absfigure}

\begin{absfigure*}{page=3, pos=t}
  \centering
  \includegraphics[width=\textwidth]{wide.pdf}
  \caption{Spanning figure on page 3}
  \label{fig:wide}
\end{absfigure*}

% A table on page 2, bottom of column 2
\begin{abstable}{page=2, pos=b, col=2}
  \centering
  \caption{Results}
  \label{tab:results}
  \begin{tabular}{lcc}
    \hline
    Method & Accuracy & F1 \\
    \hline
    Ours   & 94.2     & 0.91 \\
    \hline
  \end{tabular}
\end{abstable}

\begin{absfigure}{page=end}
  \centering
  \includegraphics[width=\columnwidth]{extra.pdf}
  \caption{Appended at the end}
  \label{fig:extra}
\end{absfigure}

See Figure~\ref{fig:one}, Table~\ref{tab:results},
and Figure~\ref{fig:extra}.

\lipsum[1-30]

\end{document}
\end{verbatim}

%% ============================================================
\section{Version history}
%% ============================================================

\begin{description}
\item[v0.1 (2026/02/25)] Initial release.
  \begin{itemize}
  \item \env{absfigure}, \env{absfigure*}, \env{abstable}, and \env{abstable*} environments
  \item Keys: \key{page}, \key{pos}, \key{col}
  \item \key{page=current} (default when \key{page} omitted)
  \item \key{page=end} support
  \item Validation of all keys with meaningful error messages
  \item End-of-document warnings for unplaced figures
  \end{itemize}
\end{description}

\end{document}
